\chapter{The Simulator and Its Targets}
\label{ch:targets}

\chapterorient{Before we open the machine, we should know what it makes. This
chapter is a guided tour of the six questions our simulator can answer, the
physical situations each one models, and the single number or table each one
produces. By the end you will be able to look at an engineering problem and say
which of the six analyses it calls for---and you will have met, in physical
form, every target the rest of the book learns to compute.}

\section{What an electromagnetic field simulator does}

An electromagnetic field simulator takes a description of a physical
device---its geometry, its materials, and how it is excited---and computes the
electromagnetic fields inside and around it. From those fields it extracts the
quantities an engineer actually wants: the capacitance between two conductors,
the inductance of a coil, the resonant frequencies of a cavity, the fraction of
a signal a microwave component reflects.

The device we use throughout this book is \emph{Palace}, a three-dimensional
finite-element solver for Maxwell's equations. Its internals are the subject of
the entire book; in this first chapter we stay outside the machine and look only
at its interface. From the outside, every run has the same shape
(\cref{fig:blackbox}): a \emph{configuration} goes in---a mesh of the geometry,
the material properties, the boundary excitations, and a choice of analysis---and
a \emph{physical product} comes out.

\begin{figure}[t]
\centering
\begin{tikzpicture}[node distance=14mm]
  \node[data, align=left, text width=24mm] (cfg)
    {geometry\\materials\\excitation\\analysis type};
  \node[stage, minimum width=34mm, minimum height=16mm, right=of cfg] (sim)
    {\textbf{the simulator}\\[1pt]\footnotesize assemble $\to$ solve $\to$ reduce};
  \node[product, align=left, text width=24mm, right=of sim] (out)
    {capacitance,\\inductance,\\frequencies,\\S-parameters,\\\dots};
  \draw[flow] (cfg) -- node[above,font=\scriptsize]{config} (sim);
  \draw[flow] (sim) -- node[above,font=\scriptsize]{product} (out);
\end{tikzpicture}
\caption[The simulator as a black box]{From the outside, every analysis is the
same black box: a configuration enters on the left and a physical product leaves
on the right. The three-word caption inside the box---\emph{assemble, solve,
reduce}---is the skeleton we will open up in \cref{ch:situating} and spend the
rest of the book filling in.}
\label{fig:blackbox}
\end{figure}

What changes from run to run is the \emph{analysis type}---the choice of which
physical question to ask. Palace offers six. They look, at first, like six
different programs. A central claim of this book, made precise in
\cref{ch:situating} and proved component by component thereafter, is that they
are really one program with a few parts swapped. But to make that claim land, we
first need the six in their own right. We take them in turn, each with the
physical picture it models and the quantity it yields.

\section{The six analyses}

\Cref{fig:sixanalyses} shows all six at a glance; the subsections that follow
walk through them one by one. Keep one eye on the figure as you read---the point
of the section is that these six very different-looking physical situations will,
by the end of the book, turn out to be six settings of the same dial.

\begin{figure}[p]
\centering
\setlength{\tabcolsep}{6pt}
\renewcommand{\arraystretch}{1.2}
\begin{tabular}{@{}cc@{}}
% --- electrostatic: capacitor ---
\begin{tikzpicture}[scale=0.95]
  \fill[simBlue!12] (-1.1,-0.7) rectangle (1.1,0.7);
  \draw[very thick,simInk] (-1.1,0.7)--(1.1,0.7);
  \draw[very thick,simInk] (-1.1,-0.7)--(1.1,-0.7);
  \node[font=\scriptsize] at (-1.4,0.7) {$+V$};
  \node[font=\scriptsize] at (-1.4,-0.7) {$0$};
  \foreach \x in {-0.8,-0.4,0,0.4,0.8}{\draw[flow,simRust,shorten >=1pt] (\x,0.62)--(\x,-0.62);}
  \node[font=\scriptsize\itshape] at (0.7,0.0) {$\vE$};
\end{tikzpicture}
&
% --- magnetostatic: coil ---
\begin{tikzpicture}[scale=0.95]
  \draw[very thick,simRust] (0,0) ellipse (0.45 and 0.9);
  \draw[-{Stealth[length=2mm]},simRust,thick] (0.45,0.2)--(0.45,-0.1);
  \node[font=\scriptsize] at (0.85,-0.4) {$I$};
  \foreach \y in {-0.5,0,0.5}{\draw[flow,simBlue] (-1.1,\y) .. controls (-0.2,\y+0.15) and (0.2,\y+0.15) .. (1.1,\y);}
  \node[font=\scriptsize\itshape] at (1.0,0.75) {$\vB$};
\end{tikzpicture}
\\[-1mm]
{\footnotesize\textbf{Electrostatic} $\to$ capacitance}
&
{\footnotesize\textbf{Magnetostatic} $\to$ inductance}
\\[3mm]
% --- eigenmode: cavity ---
\begin{tikzpicture}[scale=0.95]
  \draw[very thick,simInk] (-1.1,-0.7) rectangle (1.1,0.7);
  \draw[simTeal,thick,domain=-1.05:1.05,samples=60] plot (\x,{0.45*sin(180*\x+90)});
  \draw[simTeal!55,thick,domain=-1.05:1.05,samples=60] plot (\x,{-0.45*sin(180*\x+90)});
  \node[font=\scriptsize] at (0.0,-1.05) {$f,\,Q$};
\end{tikzpicture}
&
% --- driven: 2-port network ---
\begin{tikzpicture}[scale=0.95]
  \node[draw,thick,simInk,fill=simBgGray,minimum width=12mm,minimum height=12mm] (b) at (0,0) {\scriptsize DUT};
  \draw[flow,simGreen] (-1.5,0.25)--(b.west|-0,0.25) coordinate[midway](in);
  \draw[-{Stealth[length=2mm]},simRust] (b.west|-0,-0.25)--(-1.5,-0.25);
  \draw[flow,simGreen] (b.east)--(1.5,0);
  \node[font=\scriptsize] at (-1.2,0.55) {in};
  \node[font=\scriptsize] at (-1.25,-0.55) {refl.};
  \node[font=\scriptsize] at (1.05,-0.35) {$S_{ij}$};
\end{tikzpicture}
\\[-1mm]
{\footnotesize\textbf{Eigenmode} $\to$ frequencies \& $Q$}
&
{\footnotesize\textbf{Driven} $\to$ S-parameters}
\\[3mm]
% --- transient: pulse over time ---
\begin{tikzpicture}[scale=0.95]
  \draw[->,simGray] (-1.2,-0.55)--(1.3,-0.55) node[right,font=\scriptsize]{$t$};
  \draw[simRust,thick,domain=-1.1:1.1,samples=60]
    plot (\x,{0.7*exp(-6*\x*\x)*cos(360*\x)-0.0});
  \node[font=\scriptsize\itshape] at (-1.0,0.55) {$\vJ(t)$};
\end{tikzpicture}
&
% --- waveguide cross-section ---
\begin{tikzpicture}[scale=0.95]
  \draw[very thick,simInk] (-0.9,-0.7) rectangle (0.9,0.7);
  \foreach \x in {-0.55,-0.18,0.18,0.55}{
    \draw[-{Stealth[length=1.6mm]},simViolet] (\x,-0.5)--(\x,0.5);}
  \draw[-{Stealth[length=2.4mm]},simBlue,thick] (1.05,-0.5)--(1.05,0.5)
    node[right,font=\scriptsize]{$k_n$};
\end{tikzpicture}
\\[-1mm]
{\footnotesize\textbf{Transient} $\to$ time-domain fields}
&
{\footnotesize\textbf{Boundary mode} $\to$ waveguide modes}
\end{tabular}
\caption[The six analyses]{The six physical situations the simulator models.
Each is a different question about the same Maxwell physics; each yields a
different product. The rest of Part~I shows that all six share one skeleton, and
\cref{part:modalities} studies each in depth.}
\label{fig:sixanalyses}
\end{figure}

\subsection{Electrostatic: capacitance}

Hold two conductors at fixed voltages with an insulator (a \emph{dielectric})
between them. Charge collects on their surfaces, and an electric field $\vE$
fills the gap. The ratio of stored charge to applied voltage is the
\emph{capacitance}, the quantity that governs how a circuit stores energy in its
electric field and how quickly signals can charge a node.

The electrostatic analysis computes the capacitance of a system of conductors.
It does so by solving for the electric potential $V$ when each conductor in turn
is held at unit voltage and the others at zero, then reading the capacitance off
the resulting fields. The governing equation is the simplest in the book---a
single scalar Poisson equation, $-\Div(\varepsilon\Grad V)=0$, the same
equation that describes steady heat flow or an ideal fluid's potential. The
product is an $n\times n$ \emph{capacitance matrix} relating the charges on $n$
conductors to their voltages.

\subsection{Magnetostatic: inductance}

Run a steady current around a loop of wire and a magnetic field $\vB$ threads
through it. The ratio of magnetic flux to current is the \emph{inductance}, the
magnetic counterpart of capacitance and the quantity that governs how a circuit
stores energy in its magnetic field and resists changes in current.

The magnetostatic analysis computes inductance the same way electrostatics
computes capacitance: drive each current source in turn, solve for the field,
and reduce. The equation is the \emph{curl--curl} equation
$\Curl(\nu\,\Curl\vA)=\vJ$ for a magnetic vector potential $\vA$, where
$\nu=1/\mu$ is the reluctivity. We will see in \cref{ch:magnetostatics} that
this analysis is, almost line for line, the electrostatic analysis with the
function space changed and one weight in the final reduction swapped---our first
concrete instance of ``the same machine, one part different.''

\subsection{Eigenmode: resonant frequencies and quality factors}

Strike a bell and it rings at particular pitches; close off a region of space
with conducting walls and the electromagnetic field inside will likewise
``ring'' at particular frequencies, its \emph{resonant modes}. Each mode has a
frequency $f$ and a \emph{quality factor} $Q$ measuring how slowly it loses
energy. Cavity resonators, superconducting qubits, and filter design all live
or die by these numbers.

The eigenmode analysis finds the resonances with no external drive at all. It
assembles the operators once and hands them to a specialized eigenvalue solver,
which returns the modes as eigenpairs: each eigenvalue gives a frequency, each
eigenvector a mode field. Mathematically this is a \emph{generalized eigenvalue
problem} $\mat{K}\vx=\lambda\mat{M}\vx$. It is the one analysis whose ``solve''
is a single call to an opaque library routine rather than a loop we write
ourselves---a distinction that will turn out to matter a great deal.

\subsection{Driven: scattering parameters}

Feed a sinusoidal signal into one port of a microwave component and some of it
passes through, some reflects, and the proportions depend on frequency. The
\emph{scattering parameters}, or S-parameters, are the complex numbers $S_{ij}$
that record how much of a wave entering port $j$ leaves port $i$. They are the
universal currency of microwave and RF engineering; every filter, coupler, and
antenna is specified by its S-parameters across a frequency band.

The driven analysis computes S-parameters by solving the time-harmonic Maxwell
system at each frequency in a swept band and projecting the resulting fields onto
the ports. The operator now depends on the frequency $\omega$:
$\mat{A}(\omega)=\mat{K}+i\omega\mat{C}-\omega^2\mat{M}$. Because the operator
changes with every frequency, this analysis rebuilds and re-solves inside the
sweep---the opposite of electrostatics, which assembles once and reuses. That
contrast is the heart of \cref{ch:driven}.

\subsection{Transient: time-domain fields}

Inject a pulse of current and watch the fields evolve, step by step, in time.
The transient analysis marches the time-domain Maxwell equations forward from an
initial state, producing the full history of the fields: a movie rather than a
snapshot. It is the analysis of choice when the excitation is broadband, when
nonlinear effects matter, or when one simply wants to watch a wave propagate.

The governing equation is now a second-order system in time,
$\mat{M}\,\ddot{\vs}+\mat{C}\,\dot{\vs}+\mat{K}\,\vs=\vJ(t)$, integrated by an
ordinary-differential-equation solver. Where the driven analysis solved
independent problems at independent frequencies, the transient analysis solves a
\emph{chain}: each time step begins from the state the previous step left
behind. This is our first genuinely sequential computation, and it will be the
motivating example for the ``fold'' in \cref{ch:fold}.

\subsection{Boundary mode: waveguide propagation modes}

A waveguide---a hollow metal pipe, a coaxial line, a microstrip---carries
electromagnetic energy along its length in a set of characteristic patterns
called \emph{propagation modes}, each with a \emph{propagation constant} $k_n$
that says how fast it travels and whether it travels at all. Before the driven
analysis can excite a device through a port, it must know the modes of the
waveguide feeding that port.

The boundary-mode analysis computes those modes. It is distinguished from the
other five by a preliminary step: it extracts a two-dimensional cross-section
from the three-dimensional mesh and solves an eigenvalue problem \emph{on the
cross-section}. The solve itself is, once again, the opaque eigenvalue call of
the eigenmode analysis---now posed on a surface rather than a volume. The product
is a table of modes, each with its propagation constant and its transverse field
pattern.

\section{Six questions, one machine}

\Cref{tab:sixanalyses} lays the six side by side. Read down the columns and a
pattern jumps out that no single analysis reveals: the products differ, the
equations differ, the function spaces differ---but the \emph{structure} of every
row is identical. Each analysis assembles one or more operators from the mesh and
materials, solves something, and reduces the solution to a physical quantity.

\begin{table}[t]
\centering
\caption[The six analyses, compared]{The six analyses share a single structure:
assemble, solve, reduce. The columns that change are the physics (equation,
space) and the \emph{kind} of solve. \Cref{ch:situating} names the four kinds of
solve and the three kinds of reduction that this table is quietly tabulating.}
\label{tab:sixanalyses}
\small
\setlength{\tabcolsep}{4.5pt}
\begin{tabular}{@{}llll@{}}
\toprule
Analysis & Physical question & Governing equation & Product \\
\midrule
Electrostatic & capacitance of conductors & $-\Div(\varepsilon\Grad V)=0$ & capacitance matrix \\
Magnetostatic & inductance of current loops & $\Curl(\nu\Curl\vA)=\vJ$ & inductance matrix \\
Eigenmode & resonances of a cavity & $\mat{K}\vx=\lambda\mat{M}\vx$ & frequencies $f$, factors $Q$ \\
Driven & frequency response & $\mat{A}(\omega)\vE=\vb(\omega)$ & S-parameters $S_{ij}(\omega)$ \\
Transient & time-domain response & $\mat{M}\ddot{\vs}+\mat{C}\dot{\vs}+\mat{K}\vs=\vJ(t)$ & field history \\
Boundary mode & waveguide modes & $\mat{A}(\omega)\vx=\sigma\mat{B}\vx$ (2-D) & modes $k_n$, fields \\
\bottomrule
\end{tabular}
\end{table}

\begin{keyidea}
The simulator offers six analyses, but they are six settings of one machine.
Every one of them \emph{assembles} operators from the geometry and materials,
\emph{solves} a problem built from those operators, and \emph{reduces} the
solution to a physical quantity. What distinguishes the analyses is which
equation they discretize, on which function space, and---above all---\emph{what
kind of solve} the middle stage performs. There are only four kinds, and naming
them is the business of \cref{ch:situating}.
\end{keyidea}

This is the thesis the book exists to develop. We have it now only as a
promise---a table whose last-but-one column we have not yet learned to read. The
next chapter opens the machine far enough to see the stages \emph{assemble,
solve, reduce} as actual steps in an actual program: the lifecycle that every
one of these six analyses runs through, from the moment the user presses
\textsf{enter} to the moment the product is written out.

\summarysection
A field simulator turns a configuration---geometry, materials, excitation,
analysis type---into a physical product. Palace offers six analyses:
electrostatic (capacitance), magnetostatic (inductance), eigenmode (resonant
frequencies and quality factors), driven (S-parameters), transient (time-domain
fields), and boundary mode (waveguide propagation modes). Despite their physical
variety, all six share one skeleton: \emph{assemble, solve, reduce}. The book's
central project is to show that the variety lives almost entirely in the
``solve'' stage, of which there are only four kinds.

\furtherreading
For the physics behind the six analyses, any undergraduate electromagnetics
text---Griffiths~\citep{griffiths2017} for the fields, Pozar~\citep{pozar2011}
for the microwave quantities (S-parameters, waveguide modes, $Q$)---is ideal
background. The simulator itself is documented at the Palace
project~\citep{palace2023}. No functional-programming or finite-element
background is needed yet; we build both from scratch beginning in
\cref{part:fieldtheory} and \cref{part:framework}.

\exercisessection
\begin{enumerate}[nosep]
  \item For each of the six analyses, name one engineering device whose design
  would require it. (Example: the eigenmode analysis is essential to designing a
  superconducting-qubit resonator.)
  \item The electrostatic and magnetostatic rows of \cref{tab:sixanalyses} look
  the most alike. List every column in which they differ. (We will return to
  exactly this comparison in \cref{ch:magnetostatics}.)
  \item Which two analyses produce a result that depends on frequency? Which one
  produces a result that depends on time? What does this suggest about how their
  ``solve'' stages must differ?
  \item The driven and boundary-mode analyses both involve frequency, but only
  one of them ``rings'' with no input. Which, and why does that make its solve a
  search for eigenvalues rather than a response to a drive?
\end{enumerate}
